\chapter{Samandrag}

\subsection*{Del~\ref{part:neutron-stars} (prosjektoppgåve): Nøytronstjerner}

Nøytronstjerner inneheld materie ved ekstremt høg tettleik, der både generell
relativitetsteori og kvantefysikk trengst for å skildre strukturen. I denne
delen studerer vi korleis ulike modellar for tett materie påverkar samanhengen
mellom massen og radiusen til ei nøytronstjerne. Analysen byggjer på
Tolman--Oppenheimer--Volkoff-likningane, som blir løyste numerisk for ulike
tilstandslikningar.
Vi startar med ein enkel modell: ein ideell, ikkje-vekselverkande Fermi-gass.
Denne viser korleis degenerasjonstrykk kan motverke gravitasjonskollaps, men
òg kvifor slike idealiserte modellar ikkje er nok til å forklare observerte
nøytronstjerner.
Deretter utviklar vi ein vekselverkande modell for nukleon og lepton i likevekt.
Vekselverknadene blir skildra med sjølvkonsistente mesonfelt, som gir ei meir
realistisk tilstandslikning. Denne blir kombinert med ein standard
skorpemodell før ho blir brukt i dei numeriske stjerneutrekningane.
Dei resulterande masse--radius-relasjonane gir radiusar som er typiske for
nøytronstjerner, men dei støttar ikkje dei tyngste observerte stjernene og
skil seg tydeleg frå noverande astronomiske avgrensingar. Dette viser at den
vekselverkande modellen er ei klar forbetring frå den ideelle Fermi-gassen, men
framleis er ufullstendig ved dei høgaste tettleikane.
Samla gir denne delen ein veg frå enkle idealmodellar til ei meir realistisk
skildring av vekselverkande tett materie, og viser kvifor meir avanserte
tilstandslikningar trengst.

\subsection*{Del~\ref{part:quark-stars} (masteroppgåve): Kvarkstjerner}

Denne avhandlinga studerer kompakte stjerner frå effektive to-flavor-modellar
for kvarkmaterie. Hovudvekta ligg på QMD-modellen, der dikvark-paring gir
opphav til den farge-superleiande 2SC-fasen. Tilstandslikningane blir rekna ut
i middelfeltapproksimasjonen og brukte som input til
Tolman--Oppenheimer--Volkoff-likningane. Først blir QM-modellen brukt som eit
utgangspunkt. Her ser vi korleis sigmamassen påverkar både stivleiken til
tilstandslikninga og forma på masse--radius-kurvene til dei ulike stjernene.
Lågare sigmamassar gir større maksimale massar og sjølvbunden-liknande
oppførsel ved låge massar, medan \(m_\sigma=600~\mathrm{MeV}\) gir ei
løysing bunden av gravitasjon. Eit Bodmer--Witten-inspirert vakuumskift
reduserer masse- og radiusskalaen, men endrar ikkje denne kvalitative ordninga
i parameterområdet som er undersøkt. Hovudresultatet kjem frå QMD-modellen med
2SC-paring. Her viser vi at det fulle ein-sløyfe-potensialet må behaldast for å
få ein robust 2SC-fase. Dersom \(\Omega_{1,\mathrm{num}}\) blir fjerna, kan den
para fasen bli sterkt redusert eller forsvinne heilt. For nøytral
stjernematerie gir QMD-basismodellen ei løysing bunden av gravitasjon, med
\(M_{\max}=1.970\,M_\odot\) og \(R(M_{\max})=12.16\,\mathrm{km}\).
Parametervariasjonane viser at \(g_\Delta\), \(m_\Delta\) og \(\lambda_3\)
påverkar ulike delar av masse--radius-relasjonen. Nokre parameterendringar
aukar maksimalmassen, medan andre først og fremst reduserer radiusen eller
endrar forma på masse--radius-kurva. Utvalde kombinasjonar av parametrar gir massar
over \(2\,M_\odot\) og samsvarar med NICER-områda for både J0030+0451
og J0740+6620. Resultata tyder på at to-flavor QMD-materie med 2SC-paring kan
gi fenomenologisk lovande kompakte stjerner. Modellane som er studerte her, er
likevel reine kvarkstjernemodellar utan hadronisk skorpe, særkvarkar,
tidevassdeformabilitet eller statistisk inferens.